\chapter{Valeurs propres} \todo{Lecture 7}

\begin{definition}[Vecteur propre]
Soit $V$ un espace vectoriel sur $K$ et $f:V\to V$ un endomorphisme. Un vecteur $\boldsymbol{v}\in V\setminus \{ \boldsymbol{0} \}$ est un vecteur propre associe a la valeur propre $\lambda \in K$ si on a $f(\boldsymbol{v})=\lambda \boldsymbol{v}$
\end{definition}

\begin{example}
Soit $f:V\to V,\boldsymbol{v}\mapsto \boldsymbol{0}$ alors tout $\boldsymbol{v}\in V\setminus \{ \boldsymbol{0} \}$ est un vecteur propre associe au valeur propre $0\in K$. 
\end{example}
\begin{example}
Soit $V=\mathbf{Q}^{2}$ et $f:V\to V,\boldsymbol{x}\mapsto \begin{pmatrix}1 & 0 \\ 0 & 0\end{pmatrix}\boldsymbol{x}$ alors $\boldsymbol{e}_{1}=\begin{pmatrix}1  \\ 0\end{pmatrix}$ est un vecteur propre associe a $\lambda=1$, $\boldsymbol{e}_{2}=\begin{pmatrix}0 \\ 1 \end{pmatrix}$ est associe a $\lambda=0$ et $\boldsymbol{v}=\begin{pmatrix}1 \\ 1\end{pmatrix}$ n'est pas un vecteur propre.
\end{example}

\begin{lemma}
Soit $\dim V=n$ et $\mathscr{B}=\{ \boldsymbol{b}_{1},\dots,\boldsymbol{b}_{n} \}$ une base de $B$ et $f:V\to V$ un endomorphisme. $\mathscr{B}$ est composee de vecteurs propres si et seulement si $A^{f}_{\mathscr{B}}=\begin{pmatrix}\lambda_{1} & \dots & 0 \\ \vdots & \ddots & \vdots \\ 0 & \dots& \lambda _{n}\end{pmatrix}$ pour $\lambda _{i}\in K$.
\end{lemma}
\begin{proof}
On demontre le sens $\Rightarrow$. La $i$-eme colonne de $A_{\mathscr{B}}^{f}$ est obtenue en multipliant par $\boldsymbol{e}_{i}$ : $$A_{\mathscr{B}}^{f}\boldsymbol{e}_{i}=A_{\mathscr{B}}^{f}[\boldsymbol{b}_{i}]_{B}=[f(\boldsymbol{b}_{i})]_{\mathscr{B}}=[\lambda _{i}\boldsymbol{b}_{i}]_{\mathscr{B}}=\lambda _{i}[\boldsymbol{b}_{i}]_{\mathscr{B}}=\lambda _{i}\boldsymbol{e}_{i}$$
Le sens $\Leftarrow$ est demontre de maniere analogue.
\end{proof}
\begin{definition}
L'endomorphisme $f:V\to V$ est diagonalisable si $\boldsymbol{v}$ possede une base $\mathscr{B}=\{ \boldsymbol{b_{1}},\dots,\boldsymbol{b}_{n} \}$ composee de vecteurs propres.

Une matrice $A\in K^{n\times n}$ est diagonalisable si $f:K^{n}\to K^{n},\boldsymbol{x}\mapsto A\boldsymbol{x}$ est diagonalisable. 
\begin{exercise}
Si et seulement si il existe une matrice $P\in K^{n\times n}$ inversible telle que $P^{-1}AP$ est diagonale.
\end{exercise}
\end{definition}

\begin{lemma}
Un vecteur $\boldsymbol{v}\in V\setminus \{ \boldsymbol{0} \}$ est un vecteur propre associe a $\lambda \in K$ un valeur propre si et seulement si $\boldsymbol{v}\in \ker(f-\lambda \mathrm{Id})$.
\end{lemma}
\begin{definition}
Soit $\lambda \in K$ un valeur propre . On definit 
$$
E_{\lambda}:= \{ \boldsymbol{v}\in V:\boldsymbol{v}\in \ker(f-\lambda \mathrm{Id}) \}
$$
l'espace propre de $\lambda$. La dimension de $E_{\lambda}$ est appelle la multiplicite geometrique de $\lambda$.
\end{definition}

\begin{theorem}
Soient $\boldsymbol{v}_{1},\dots,\boldsymbol{v}_{k}\in V$ des vecteurs propres associe aux valeurs propres $\lambda_{1},\dots,\lambda _{k}\in K$ distincts. Alors $\{ \boldsymbol{v}_{1},\dots,\boldsymbol{v}_{k} \}$ est un ensemble libre.
\end{theorem}
\begin{proof}
Pour le cas $k=2$ on suppose que l'ensemble n'est pas libre. Alors ils existent $\alpha_{1},\alpha_{2}\in K\setminus \{ 0 \}$ tels que $$\alpha_{1}\boldsymbol{v}_{1}+\alpha_{2}\boldsymbol{v}_{2}=\boldsymbol{0}=f(\boldsymbol{0})=\alpha_{1}\lambda_{1}\boldsymbol{v}_{1}+\alpha_{2}\lambda_{2}\boldsymbol{v}_{2}$$Sans perte de generalite on pose que $\lambda_{2}\neq 0$, alors

\begin{align*}
\frac{\alpha_{1}\lambda_{1}}{\lambda_{2}}\boldsymbol{v}_{1}+\alpha_{2}\boldsymbol{v}_{2} & =0 \quad |-\alpha_{1}\boldsymbol{v}_{1}+\alpha_{2}\boldsymbol{v}_{2} \\
\underbrace{ \alpha_{1} }_{ \neq 0 }\underbrace{ \left( \frac{\lambda_{1}}{\lambda_{2}}-1 \right) }_{ \neq 0 }\boldsymbol{v}_{1} & =0 
\end{align*}

qui pose une contradiction car les vecteurs propres ne sont jamais nuls.

Cas general: Ils existent $\alpha_{1},\dots,\alpha _{k}\in K\setminus \{ 0 \}$ tels que 

\begin{align}
\alpha_{1}\boldsymbol{v}_{1}+\dots+\alpha _{k}\boldsymbol{v}_{k} & =\boldsymbol{0} \label{temp23}\\
\frac{\alpha_{1}\lambda_{1}}{\lambda _{k}}\boldsymbol{v}_{1}+..+\alpha _{k}\boldsymbol{v}_{k} & =\boldsymbol{0}  & |-\ref{temp23}\\
\alpha_{1}\left( \frac{\lambda_{1}}{\lambda _{k}}-1 \right)\boldsymbol{v}_{1}+\dots+\alpha _{k-1}\left( \frac{\lambda_{k-1}}{\lambda _{k}}-1 \right)\boldsymbol{v}_{k-1}=\boldsymbol{0}
\end{align}


\end{proof}


\begin{example}
Soit $A=\begin{pmatrix}0 & 1 \\ 0 & 0\end{pmatrix}\in K^{2\times2}$. On cherche $\ker(A-\lambda I)\supsetneq \{ 0 \}$ et on trouve que $\lambda=0$ est le seul valeur propre de $A$.


Si $\mathscr{B}=\{v_{1},v_{2}\}\subseteq K^{2}$ est une base de vecteurs propres de $A$, alors tout $\boldsymbol{x}\in K^{2}$ tel que $\boldsymbol{x}=\alpha v_{1}+\beta v_{2}$ on a que $A\boldsymbol{x}=\alpha Av_{1}+\beta Av_{2}=0=0$ mais $A\boldsymbol{e}_{2}\neq 0$. Alors $A$ n'est pas diagonalisable.
\end{example}
\begin{example}
$A=\begin{pmatrix}1 & 1 \\ 0 & 1\end{pmatrix}\in K^{2\times 2}$. On a $\lambda=1$ le seul valeur propre. Si $\mathscr{B}=\{b_{1},b_{2} \}\subseteq K^{2}$ est une base de vecteurs propres alors pour $[x]_{\mathscr{B}}=\begin{pmatrix}\alpha \\ \beta\end{pmatrix}\in K^{2\times1}$ alors
\begin{align*}
A\boldsymbol{x} & =\alpha Ab_{1}+\beta Ab_{2} \\
 & =\alpha b_{1}+\beta b_{2} \\
 & =\boldsymbol{x}
\end{align*}
mais $A\boldsymbol{e}_{2}\neq \boldsymbol{e}_{2}$.
\end{example}
\begin{example}
Soit $A=\begin{pmatrix}\lambda & 1 \\ 0 & \lambda \end{pmatrix}$. $\lambda$ est le seul valeur propre de $A$. S'il existe une base de vecteurs propres de $A$ alors $\forall \boldsymbol{x}\in K^{2}:A\boldsymbol{x}=\boldsymbol{x}\boldsymbol{x}$ mais $A\boldsymbol{e}_{2}=\begin{pmatrix}1 \\ \lambda\end{pmatrix}\neq \lambda \begin{pmatrix}0 \\ 1\end{pmatrix}$
\end{example}

\begin{theorem}
Soit $f:V\to V$ un endomorphisme et $\lambda_{1},\dots,\lambda _{k}\in K$ les valeurs propres distincts de $f$. Soit $g_{i}=\dim E_{\lambda _{i}}$ pour $i=1,\dots ,k$, $f$ est diagonalisable si et seulement si $\sum_{i=1}^{k}g_{i}=n$.
\end{theorem}
\begin{proof}
\begin{itemize}
\item $\Leftarrow$: Soient $\{ v_{1}^{i},\dots,v_{g_{i}}^{i} \}$ une base de $E_{\lambda _{i}}$ pour $i=1,\dots,k$. On a que
\begin{equation}
v_{1}^{1},\dots,v_{g_{1}}^{1},v_{1}^{2},\dots,v_{g_{2}}^{2},\dots,v_{1}^{k},\dots,v_{g_{k}}^{k} \label{baseV}
\end{equation}
est libre, parce que avec $\alpha _{i}^{j}\in K$ pour $i=1,\dots,k$  et $1\le j \le g_{i}$ et 
$$
\sum_{i=1}^{k} \sum_{j=1}^{g_{i}} \alpha _{i}^{j}v_{j}^{i}=0
$$
parce que les vecteurs propres associes aux valeurs propres distincts sont libres donc $\sum_{i=1}^{g_{i}}\alpha _{i}^{j}v_{j}^{i}=0$ et tous $\alpha _{i}^{j}=0$.

Alors, si $g_{1}+\dots +g_{k}=n$ on a que \ref{baseV} est une base de $V$ composee de vecteurs propres donc $f$ est diagonalisable.
\item $\Rightarrow$: On vient de montrer que $\sum _{i=1}^{k} g_{i}\le n$. Soit
$$
v_{1}^{1},\dots,v_{u_{1}}^{1},v_{1}^{2},\dots,v_{u_{2}}^{2},\dots,v_{1}^{k},\dots,v_{u_{k}}^{k}
$$
une base de $V$ composee de vecteurs propres. Alors :
\begin{enumerate}
\item $u_{1}+\dots+u_{k}=n$
\item Mais $u_{i}\le g_{i}$.
\end{enumerate}
\end{itemize}
donc on a 
$$
n=u_{1}+\dots+u_{k}\le g_{1}+\dots +g_{k}\le n
$$
alors par double inclusion $g_{1}+\dots g_{k}=n$.
\end{proof}


\subsection*{Comment decider si \texorpdfstring{$f$}{f} est diagonalisable}

\begin{enumerate}
\item Trouver tous $\lambda_{1},\dots,\lambda _{k}\in K$ valeurs propres distincts.
\item Calculer $g_{i}=\dim(\ker(f-\lambda _{i}\mathrm{Id}))$ par Gauss.
\item Si $\sum_{i=1}^{k}g_{i}=n$ alors $f$ est diagonalisable sinon non.
\end{enumerate}

\begin{exercise}[Exo etoile serie 3]
Soit $K$ un corps et $f,g\in K[x]$ deux polynomes pas tous les deux nuls. Considerons l'ensemble des diviseurs communs a $f$ et $g$:
$$
D_{f,g}=\{ d\in K[x]:d\mid f , d\mid g\}
$$
\begin{enumerate}
\item Montrer qu'il existe un unique polynome $d\in D_{f,g}$ unitaire et de degre maximal.
\item Montrer que $d=\gcd(f,g)$.
\end{enumerate}
\end{exercise}
\begin{proof}[Solution]
\begin{enumerate}
\item Sans perte de generalite, on suppose que $g\neq 0$. Grace a la division euclidienne on peut ecrire $f=qg+r$ avec $q,r\in K[x]$ et $\deg r< \deg g$, on remarqiue ainsi que $D_{f,g}=D_{g,r}$, pourquoi? En effet soit $h\in K[x]$ alors $h \mid f$ et $h \mid g$ puis $h \mid g$ et $h \mid r$. En continuant ainsi (par des divisions euclidiennes succesives), on arrive a l'existence de $r^{*}\in K[x]$ tel que $D_{f,g}=D_{r^{*},0}$. Cette procedure se termine car la somme de degres de polynomes est strictement decroissante. On a
$$
D_{f,g}=D_{r^{*},0}=\{ d\in K[x]:d\mid r^{*} \}
$$
et le degre des polynomes dans $D_{f,g}$ est borne par $\deg r^{*}$. Soit $a_{n}\in K$ le coefficient dominant de $r^{*}$, on pose
$$
d^{*}:= \frac{1}{a_{n}}r^{*}
$$
et on observe que $d^{*}$ est unitaire et de degre maximal par construction.

Pour monter l'unicite, on prend $\tilde{d}\in D_{f,g}$ de degre maximal et unitaire alors $\tilde{d}\in D_{r^{*},0}=D_{d^{*},0}$ qui implique qu'il existe $\lambda \in K$ tel que $\lambda \tilde{d}=d^{*}$ (car $\deg \tilde{d} = \deg d^{*}$). En comparant les coefficients dominants de $\lambda \tilde{d}$ et $d^{*}$ on conclut que $\lambda=1$ et donc $\tilde{d}=d^{*}$ qui demontre l'unicite de $d^{*}$.
\item Il suffit de montrer que tout diviseur commun de $f$ et $g$ divise $d^{*}$, car $D_{f,g}=D_{d^{*},0}$, et que $d^{*}$ divise $f$ et $g$ car $d^{*}\in D_{f,g}$. On conclut par l'unicite de $d^{*}$.
\end{enumerate}
\end{proof}
\begin{exercise}[Exo plus serie 4]
Soit $V$ un espace vectoriel sur un corps $K$ de dimension finie et $f:V\to V$ un endomorphisme. Soit $p:V\to V$ un automorphisme (c-a-d un endomorphisme inversible). Montrer que $\lambda \in K$ est une valeur propre de $f$ si et seulement si $\lambda$ est une valeur propre de l'endomorphisme $p ^{-1}\circ f \circ p$.
\end{exercise}
\begin{proof}[Solution]
\begin{itemize}
\item $\Rightarrow$: Supposons que $\lambda \in K$ est une valeur propre de $f$ avec $\boldsymbol{v}\in V\setminus \{ \boldsymbol{0} \}$ un vectuer propre associe. Posons $\boldsymbol{w}:=p ^{-1}(\boldsymbol{v})\neq \boldsymbol{0}$ (non nul car $p$ est un automorphisme). Si on applique la composition
\begin{align*}
(p ^{-1} \circ f \circ p)(\boldsymbol{w}) & =(p ^{-1} \circ f \circ p)(p ^{-1}(\boldsymbol{v})) \\
 & =p ^{-1} (f(\boldsymbol{v})) \\
 & = p ^{-1} (\lambda \boldsymbol{v}) \\
 & =\lambda p ^{-1} (\boldsymbol{v}) \\
 & =\lambda \boldsymbol{w}
\end{align*}

\item $\Leftarrow$: Supposons que $\lambda \in K$ est un valeur propre de $p ^{-1} \circ f \circ p$ et $\boldsymbol{w}\in V\setminus \{ \boldsymbol{0} \}$ est un vecteur propre associe. Posons $\boldsymbol{v}=p(\boldsymbol{w})\neq 0$ donc
\begin{align*}
f(\boldsymbol{v}) & =f(p(\boldsymbol{w})) \\
 & =(p \circ p ^{-1})(f(p(\boldsymbol{w}))) \\
 & =p(p ^{-1} \circ f \circ p)(\boldsymbol{w}) \\
 & = p(\lambda \boldsymbol{w}) \\  & =\lambda p(\boldsymbol{w})\\
 & = \lambda \boldsymbol{v}
\end{align*}
\end{itemize}
\end{proof}